% conclusion.tex -- EVM Executor Paper (BNR-2026-008)

\section{Conclusion}

This paper investigated the feasibility of using a minimal go-ethereum fork as
the EVM execution engine for the Basis Network enterprise zkEVM L2. Through
systematic analysis of five production L2 systems, a complete mapping of 75
Cancun EVM opcodes to ZK constraint costs, and the design of a selective ZK
tracer, we establish the following conclusions:

\begin{enumerate}
  \item The execution-proving separation pattern is universal across production
    zkEVM systems. Geth's role is transaction execution and trace generation;
    a separate prover (in our case, a Rust-based PLONK prover communicating
    via gRPC) consumes these traces for proof generation.

  \item A minimal Geth integration via Go module import (Strategy~A) is the
    recommended approach, requiring approximately 24{,}500~LOC of Geth's core
    packages with zero lines of EVM modification. This is validated by
    Optimism's op-geth achieving a complete L2 with only 34~EVM lines changed,
    and Scroll's subsequent migration away from tight fork coupling.

  \item The \texttt{core/tracing} hooks API provides an efficient mechanism
    for ZK-specific trace capture, with projected overhead of 10--30\% for
    selective tracing versus 10--50$\times$ for full opcode capture.

  \item Projected throughput of 4{,}000--12{,}000~tx/s for simple transfers
    and 1{,}500--6{,}000~tx/s for storage writes with ZK tracing comfortably
    exceeds the 1{,}000~tx/s hypothesis threshold.

  \item KECCAK256 (${\sim}$150{,}000 R1CS constraints) and storage operations
    (255 Poseidon operations per SLOAD/SSTORE) dominate ZK proving costs.
    KECCAK256 mitigation via preimage oracle with lookup tables is a mandatory
    circuit-level optimization.

  \item Six architectural invariants (INV-EVM-1 through INV-EVM-6) are
    established as constraints for downstream pipeline stages: the Logicist
    (TLA+ specification), the Architect (Go/Rust implementation), and the
    Prover (Coq verification).
\end{enumerate}

\subsection{Recommendations}

For the Basis Network enterprise zkEVM L2 implementation:

\begin{enumerate}
  \item Import go-ethereum v1.14+ as a Go module dependency. Use
    \texttt{core/vm}, \texttt{core/state}, \texttt{core/types},
    \texttt{core/tracing}, and supporting packages.

  \item Implement a custom \texttt{StateDB} satisfying the \texttt{vm.StateDB}
    interface, backed by a Poseidon-based sparse Merkle tree for ZK-friendly
    state management.

  \item Deploy the selective ZK tracer (state changes only) in production, with
    full opcode capture available as a debug mode.

  \item Target Geth v1.14.12+ for the Cancun EVM feature set. Do not set
    PragueTime in the chain configuration.

  \item Monitor Scroll's Euclid migration and Geth upstream releases for
    tracing API changes that may affect the \texttt{core/tracing.Hooks}
    interface.
\end{enumerate}

\subsection{Future Work}

Immediate next steps in the Basis Network R\&D pipeline:

\begin{itemize}
  \item \textbf{RU-L2 (State Management)}: Design and benchmark the Poseidon
    SMT implementation behind the \texttt{vm.StateDB} interface.
  \item \textbf{RU-L3 (Witness Generation)}: Specify the exact trace format
    consumed by the Rust prover and validate end-to-end trace completeness.
  \item \textbf{Actual benchmarks}: Execute the Go benchmark suite on target
    hardware to validate projected throughput figures.
  \item \textbf{KECCAK256 circuit}: Design and benchmark the preimage oracle
    with lookup tables for KECCAK256 proving.
\end{itemize}
